problem:  
Let \((M^n,g)\) be a compact Riemannian manifold with smooth, umbilic, but not totally geodesic boundary \(\partial M\), and \(n \ge 3\).  
Consider the conformal Laplacian and boundary conformal operator given by  
\[
L_g = -a_n \Delta_g + R_g, \quad B_g = \frac{\partial}{\partial \nu_g} + b_n H_g,
\]
where \(a_n = \frac{4(n-1)}{n-2}\) and \(b_n = \frac{2}{n-2}\), with \(R_g\) the scalar curvature and \(H_g\) the mean curvature of \(\partial M\).  

Let \(u>0\) be a smooth function, and let \(\tilde g = u^{4/(n-2)} g\). Then  
\[
L_g u = c u^{\frac{n+2}{n-2}} \text{ in } M, \quad B_g u = H_0 u^{\frac{n}{n-2}} \text{ on } \partial M,
\]
where \(c\) is a constant.  

It is known (Escobar, Marques, Almaraz, et al.) that when \(H_0\) is constant, existence of such \(u\) is largely understood. However, suppose now that \(H_0 = f\) is a *prescribed smooth, nonconstant boundary function* satisfying  
\[
\int_{\partial M} (\partial_\nu f + \lambda f) \, dS_g = 0
\]
for all conformal Killing fields \(\nu\) on \((M,g)\) (a Kazdan–Warner–type constraint).  

**Open problem:** Determine whether the above condition (or any natural finite set of conformal integral constraints) is *sufficient* for the existence of a positive solution \(u\) to the boundary Yamabe equation system, ensuring that \(\tilde g = u^{4/(n-2)}g\) has constant scalar curvature and prescribed mean curvature \(H_{\tilde g}=f\) on \(\partial M\), at least when the conformal class \([g]\) has positive Yamabe invariant \(Y(M,\partial M,[g])>0.\)  

Equivalently, is there a Kazdan–Warner–type obstruction that completely characterizes solvability for the boundary mean curvature prescription problem under the conformal deformation preserving constant scalar curvature?  

Why is it a "good" problem:  
This question refines the boundary Yamabe problem into a precise, conjectural framework seeking necessary and sufficient Kazdan–Warner–type obstructions for prescribing general boundary mean curvature functions. It is both simple and profound: it asks whether the integral invariants arising from conformal symmetries fully govern the solvability of a highly nonlinear conformally invariant elliptic system. It sits at the intersection of geometric analysis, conformal geometry, and nonlinear PDE theory. The problem’s clarity (a single integral obstruction sufficiency question) makes it sharply formulated, while its solution would unify diverse results on scalar and mean curvature prescription, extending deep structures from closed manifolds to manifolds with boundary. It thus embodies the ideals of a “good” problem—conceptually simple, geometrically rich, analytically deep, and capable of linking multiple mathematical domains.